% ── §7 Methods: Dataset ──
\section{Dataset}
\label{sec:methods-data}

\subsection{Primary Benchmark: PSICOV150}
\label{sec:psicov-data}

Our primary evaluation dataset comprises the 150 protein domains introduced
by Jones et al.\ \cite{jones2012} as a benchmark for sparse inverse covariance
estimation in contact prediction.  Each target consists of a deep multiple
sequence alignment generated by three iterations of JackHMMER against
UniRef100, together with a repaired and renumbered native experimental
structure from the Protein Data Bank.  Alignment depths range from 511 to
74{,}836 sequences (median 3{,}228), and protein widths range from 50 to 266
residues (median 140).

A critical structural property of this dataset is documented in its README:
the number of alignment columns equals exactly the number of C-$\alpha$ atom
records in each PDB file, because gap columns relative to the target sequence
were removed during alignment construction.  We verified this identity mapping
empirically for 149 of the 150 targets; the single exception (target
\texttt{1vjkA}) carries one extra C-terminal residue in the PDB file that is
absent from the target sequence, which we handle by restricting contact
computation to residues numbered within the alignment width.  This identity
mapping eliminates the need for sequence-to-structure alignment and removes an
entire class of potential errors.

\subsection{Ground Truth}
\label{sec:ground-truth}

Native contacts are defined as pairs of residues $(r_a, r_b)$ satisfying
\begin{equation}
	d(\mathrm{C}\beta_{r_a},\, \mathrm{C}\beta_{r_b}) < 8.0~\text{\AA}
	\quad \wedge \quad |r_a - r_j| \geq 6,
	\label{eq:contact-def}
\end{equation}
where $\mathrm{C}\beta$ coordinates are used for all residue types except
glycine, for which $\mathrm{C}\alpha$ is substituted.  The 8~\AA{} cutoff and
minimum separation of 6 residues follow the standard conventions established
in the CASP contact prediction assessments~\cite{ezkurdia2009}.  Across the
150 targets, we identify 47{,}430 native contacts among 1{,}675{,}886 candidate
pairs, yielding a base contact rate of 2.83\%.

As an independent literature baseline, we use the contact predictions shipped
with the PSICOV supplementary data (\texttt{con/*.out} files).  These files
enumerate candidate pairs sorted by prediction score, providing a direct
comparison point for our methods against the state-of-the-art sparse inverse
covariance approach on identical data.

\subsection{Structural Models}
For cross-protein transfer experiments requiring per-sequence structural
labels, we built homology models of all 299 unique Omicron Spike sequences
using SWISS-MODEL's automated pipeline~\cite{waterhouse2018}, obtaining
1{,}128 PDB files with GMQE scores ranging from 0.69 to 0.72.
